%% figures/fig-simultaneity.tex -- Figure for Sec. 2.1.4.
%% Pulled into notes.tex with %% figures/fig-simultaneity.tex -- Figure for Sec. 2.1.4.
%% Pulled into notes.tex with %% figures/fig-simultaneity.tex -- Figure for Sec. 2.1.4.
%% Pulled into notes.tex with %% figures/fig-simultaneity.tex -- Figure for Sec. 2.1.4.
%% Pulled into notes.tex with \input{figures/fig-simultaneity}.
%% Needs tikz with the arrows.meta library, loaded in the main preamble.
%%
%% Geometry: the world line is x = 0.18 t^2, so the velocity at time t is
%% v = 0.36 t. In an (x,t) diagram the simultaneity line of the instantaneous
%% rest frame at an event has slope dt/dx = v there, so the two slices drawn
%% in panel (a) have slopes 0.432 and 0.648 and cannot be parallel. Every
%% number below follows from those two lines; nothing is drawn by eye.

\begin{figure}[htb]
\centering
\begin{tikzpicture}[>={Stealth[length=5pt]},scale=2.1,line join=round]

%% ---------------- panel (a): the particle's own simultaneity ----------
\begin{scope}
  \fill[gray!20] (-0.5,0.872) -- (1.35,1.671) -- (1.35,2.297) -- (-0.5,1.098)
     -- cycle;
  %% the two rest-frame slices
  \draw (-0.5,0.872) -- (1.35,1.671);
  \draw (-0.5,1.098) -- (1.35,2.297);
  \node[right,font=\footnotesize] at (1.36,1.66) {$\tau$};
  \node[right,font=\footnotesize] at (1.36,2.30) {$\tau+\Delta\tau$};
  %% world line
  \draw[very thick] plot[domain=0.62:2.6,samples=60,smooth]
       ({0.18*\x*\x},\x);
  \node[left] at (0.02,2.5) {$\gamma$};
  %% acceleration
  \draw[->,thick] (0.259,1.2) -- (0.72,1.2);
  \node[below,font=\footnotesize] at (0.49,1.19) {$a$};
  %% the gap, measured at both ends
  \draw[<->] (-0.4,0.915) -- (-0.4,1.163);
  \draw[<->] (1.25,1.628) -- (1.25,2.232);
  \node[left,font=\footnotesize]  at (-0.42,1.04) {narrower};
  \node[right,font=\footnotesize] at (1.27,1.93) {wider};
  \node at (0.42,0.12) {(a) Dirac};
\end{scope}

%% ---------------- panel (b): constant time in one frame ---------------
\begin{scope}[xshift=3.25cm]
  \fill[gray!20] (-0.5,1.27) rectangle (1.35,1.694);
  \draw (-0.5,1.27)  -- (1.35,1.27);
  \draw (-0.5,1.694) -- (1.35,1.694);
  \node[right,font=\footnotesize] at (1.36,1.27)  {$t$};
  \node[right,font=\footnotesize] at (1.36,1.694) {$t+\Delta t$};
  \draw[very thick] plot[domain=0.62:2.6,samples=60,smooth]
       ({0.18*\x*\x},\x);
  \node[left] at (0.02,2.5) {$\gamma$};
  \draw[->,thick] (0.259,1.2) -- (0.72,1.2);
  \node[below,font=\footnotesize] at (0.49,1.19) {$a$};
  \draw[<->] (-0.4,1.27) -- (-0.4,1.694);
  \draw[<->] (1.25,1.27) -- (1.25,1.694);
  \node[left,font=\footnotesize]  at (-0.46,1.48) {uniform};
  \node at (0.42,0.12) {(b) Lorentz};
\end{scope}

\end{tikzpicture}
\caption{Where the $4/3$ goes. Both panels show the same accelerating world
line $\gamma$ and the same slab of it, and differ only in what is being called
``at the same moment''. \emph{(a)} Dirac's tube is built from the particle's own
succession of rest frames. Because the velocity changes along $\gamma$,
successive slices are not parallel: they fan open on the side the particle is
accelerating towards, and the slab is wider there. That fanning is the factor
$1+\varepsilon\,a\cdot n$ of Eq.~(\ref{eq:tubeelement}). \emph{(b)} Collecting
field momentum at constant time in one chosen frame gives parallel slices and a
slab of uniform thickness, with no such factor. The wedge by which the two
shaded regions differ is worth exactly one third of the electromagnetic mass,
which is the whole of the discrepancy between Eqs.~(\ref{eq:massren}) and
(\ref{eq:divterm}).
\label{fig:simultaneity}}
\end{figure}
.
%% Needs tikz with the arrows.meta library, loaded in the main preamble.
%%
%% Geometry: the world line is x = 0.18 t^2, so the velocity at time t is
%% v = 0.36 t. In an (x,t) diagram the simultaneity line of the instantaneous
%% rest frame at an event has slope dt/dx = v there, so the two slices drawn
%% in panel (a) have slopes 0.432 and 0.648 and cannot be parallel. Every
%% number below follows from those two lines; nothing is drawn by eye.

\begin{figure}[htb]
\centering
\begin{tikzpicture}[>={Stealth[length=5pt]},scale=2.1,line join=round]

%% ---------------- panel (a): the particle's own simultaneity ----------
\begin{scope}
  \fill[gray!20] (-0.5,0.872) -- (1.35,1.671) -- (1.35,2.297) -- (-0.5,1.098)
     -- cycle;
  %% the two rest-frame slices
  \draw (-0.5,0.872) -- (1.35,1.671);
  \draw (-0.5,1.098) -- (1.35,2.297);
  \node[right,font=\footnotesize] at (1.36,1.66) {$\tau$};
  \node[right,font=\footnotesize] at (1.36,2.30) {$\tau+\Delta\tau$};
  %% world line
  \draw[very thick] plot[domain=0.62:2.6,samples=60,smooth]
       ({0.18*\x*\x},\x);
  \node[left] at (0.02,2.5) {$\gamma$};
  %% acceleration
  \draw[->,thick] (0.259,1.2) -- (0.72,1.2);
  \node[below,font=\footnotesize] at (0.49,1.19) {$a$};
  %% the gap, measured at both ends
  \draw[<->] (-0.4,0.915) -- (-0.4,1.163);
  \draw[<->] (1.25,1.628) -- (1.25,2.232);
  \node[left,font=\footnotesize]  at (-0.42,1.04) {narrower};
  \node[right,font=\footnotesize] at (1.27,1.93) {wider};
  \node at (0.42,0.12) {(a) Dirac};
\end{scope}

%% ---------------- panel (b): constant time in one frame ---------------
\begin{scope}[xshift=3.25cm]
  \fill[gray!20] (-0.5,1.27) rectangle (1.35,1.694);
  \draw (-0.5,1.27)  -- (1.35,1.27);
  \draw (-0.5,1.694) -- (1.35,1.694);
  \node[right,font=\footnotesize] at (1.36,1.27)  {$t$};
  \node[right,font=\footnotesize] at (1.36,1.694) {$t+\Delta t$};
  \draw[very thick] plot[domain=0.62:2.6,samples=60,smooth]
       ({0.18*\x*\x},\x);
  \node[left] at (0.02,2.5) {$\gamma$};
  \draw[->,thick] (0.259,1.2) -- (0.72,1.2);
  \node[below,font=\footnotesize] at (0.49,1.19) {$a$};
  \draw[<->] (-0.4,1.27) -- (-0.4,1.694);
  \draw[<->] (1.25,1.27) -- (1.25,1.694);
  \node[left,font=\footnotesize]  at (-0.46,1.48) {uniform};
  \node at (0.42,0.12) {(b) Lorentz};
\end{scope}

\end{tikzpicture}
\caption{Where the $4/3$ goes. Both panels show the same accelerating world
line $\gamma$ and the same slab of it, and differ only in what is being called
``at the same moment''. \emph{(a)} Dirac's tube is built from the particle's own
succession of rest frames. Because the velocity changes along $\gamma$,
successive slices are not parallel: they fan open on the side the particle is
accelerating towards, and the slab is wider there. That fanning is the factor
$1+\varepsilon\,a\cdot n$ of Eq.~(\ref{eq:tubeelement}). \emph{(b)} Collecting
field momentum at constant time in one chosen frame gives parallel slices and a
slab of uniform thickness, with no such factor. The wedge by which the two
shaded regions differ is worth exactly one third of the electromagnetic mass,
which is the whole of the discrepancy between Eqs.~(\ref{eq:massren}) and
(\ref{eq:divterm}).
\label{fig:simultaneity}}
\end{figure}
.
%% Needs tikz with the arrows.meta library, loaded in the main preamble.
%%
%% Geometry: the world line is x = 0.18 t^2, so the velocity at time t is
%% v = 0.36 t. In an (x,t) diagram the simultaneity line of the instantaneous
%% rest frame at an event has slope dt/dx = v there, so the two slices drawn
%% in panel (a) have slopes 0.432 and 0.648 and cannot be parallel. Every
%% number below follows from those two lines; nothing is drawn by eye.

\begin{figure}[htb]
\centering
\begin{tikzpicture}[>={Stealth[length=5pt]},scale=2.1,line join=round]

%% ---------------- panel (a): the particle's own simultaneity ----------
\begin{scope}
  \fill[gray!20] (-0.5,0.872) -- (1.35,1.671) -- (1.35,2.297) -- (-0.5,1.098)
     -- cycle;
  %% the two rest-frame slices
  \draw (-0.5,0.872) -- (1.35,1.671);
  \draw (-0.5,1.098) -- (1.35,2.297);
  \node[right,font=\footnotesize] at (1.36,1.66) {$\tau$};
  \node[right,font=\footnotesize] at (1.36,2.30) {$\tau+\Delta\tau$};
  %% world line
  \draw[very thick] plot[domain=0.62:2.6,samples=60,smooth]
       ({0.18*\x*\x},\x);
  \node[left] at (0.02,2.5) {$\gamma$};
  %% acceleration
  \draw[->,thick] (0.259,1.2) -- (0.72,1.2);
  \node[below,font=\footnotesize] at (0.49,1.19) {$a$};
  %% the gap, measured at both ends
  \draw[<->] (-0.4,0.915) -- (-0.4,1.163);
  \draw[<->] (1.25,1.628) -- (1.25,2.232);
  \node[left,font=\footnotesize]  at (-0.42,1.04) {narrower};
  \node[right,font=\footnotesize] at (1.27,1.93) {wider};
  \node at (0.42,0.12) {(a) Dirac};
\end{scope}

%% ---------------- panel (b): constant time in one frame ---------------
\begin{scope}[xshift=3.25cm]
  \fill[gray!20] (-0.5,1.27) rectangle (1.35,1.694);
  \draw (-0.5,1.27)  -- (1.35,1.27);
  \draw (-0.5,1.694) -- (1.35,1.694);
  \node[right,font=\footnotesize] at (1.36,1.27)  {$t$};
  \node[right,font=\footnotesize] at (1.36,1.694) {$t+\Delta t$};
  \draw[very thick] plot[domain=0.62:2.6,samples=60,smooth]
       ({0.18*\x*\x},\x);
  \node[left] at (0.02,2.5) {$\gamma$};
  \draw[->,thick] (0.259,1.2) -- (0.72,1.2);
  \node[below,font=\footnotesize] at (0.49,1.19) {$a$};
  \draw[<->] (-0.4,1.27) -- (-0.4,1.694);
  \draw[<->] (1.25,1.27) -- (1.25,1.694);
  \node[left,font=\footnotesize]  at (-0.46,1.48) {uniform};
  \node at (0.42,0.12) {(b) Lorentz};
\end{scope}

\end{tikzpicture}
\caption{Where the $4/3$ goes. Both panels show the same accelerating world
line $\gamma$ and the same slab of it, and differ only in what is being called
``at the same moment''. \emph{(a)} Dirac's tube is built from the particle's own
succession of rest frames. Because the velocity changes along $\gamma$,
successive slices are not parallel: they fan open on the side the particle is
accelerating towards, and the slab is wider there. That fanning is the factor
$1+\varepsilon\,a\cdot n$ of Eq.~(\ref{eq:tubeelement}). \emph{(b)} Collecting
field momentum at constant time in one chosen frame gives parallel slices and a
slab of uniform thickness, with no such factor. The wedge by which the two
shaded regions differ is worth exactly one third of the electromagnetic mass,
which is the whole of the discrepancy between Eqs.~(\ref{eq:massren}) and
(\ref{eq:divterm}).
\label{fig:simultaneity}}
\end{figure}
.
%% Needs tikz with the arrows.meta library, loaded in the main preamble.
%%
%% Geometry: the world line is x = 0.18 t^2, so the velocity at time t is
%% v = 0.36 t. In an (x,t) diagram the simultaneity line of the instantaneous
%% rest frame at an event has slope dt/dx = v there, so the two slices drawn
%% in panel (a) have slopes 0.432 and 0.648 and cannot be parallel. Every
%% number below follows from those two lines; nothing is drawn by eye.

\begin{figure}[htb]
\centering
\begin{tikzpicture}[>={Stealth[length=5pt]},scale=2.1,line join=round]

%% ---------------- panel (a): the particle's own simultaneity ----------
\begin{scope}
  \fill[gray!20] (-0.5,0.872) -- (1.35,1.671) -- (1.35,2.297) -- (-0.5,1.098)
     -- cycle;
  %% the two rest-frame slices
  \draw (-0.5,0.872) -- (1.35,1.671);
  \draw (-0.5,1.098) -- (1.35,2.297);
  \node[right,font=\footnotesize] at (1.36,1.66) {$\tau$};
  \node[right,font=\footnotesize] at (1.36,2.30) {$\tau+\Delta\tau$};
  %% world line
  \draw[very thick] plot[domain=0.62:2.6,samples=60,smooth]
       ({0.18*\x*\x},\x);
  \node[left] at (0.02,2.5) {$\gamma$};
  %% acceleration
  \draw[->,thick] (0.259,1.2) -- (0.72,1.2);
  \node[below,font=\footnotesize] at (0.49,1.19) {$a$};
  %% the gap, measured at both ends
  \draw[<->] (-0.4,0.915) -- (-0.4,1.163);
  \draw[<->] (1.25,1.628) -- (1.25,2.232);
  \node[left,font=\footnotesize]  at (-0.42,1.04) {narrower};
  \node[right,font=\footnotesize] at (1.27,1.93) {wider};
  \node at (0.42,0.12) {(a) Dirac};
\end{scope}

%% ---------------- panel (b): constant time in one frame ---------------
\begin{scope}[xshift=3.25cm]
  \fill[gray!20] (-0.5,1.27) rectangle (1.35,1.694);
  \draw (-0.5,1.27)  -- (1.35,1.27);
  \draw (-0.5,1.694) -- (1.35,1.694);
  \node[right,font=\footnotesize] at (1.36,1.27)  {$t$};
  \node[right,font=\footnotesize] at (1.36,1.694) {$t+\Delta t$};
  \draw[very thick] plot[domain=0.62:2.6,samples=60,smooth]
       ({0.18*\x*\x},\x);
  \node[left] at (0.02,2.5) {$\gamma$};
  \draw[->,thick] (0.259,1.2) -- (0.72,1.2);
  \node[below,font=\footnotesize] at (0.49,1.19) {$a$};
  \draw[<->] (-0.4,1.27) -- (-0.4,1.694);
  \draw[<->] (1.25,1.27) -- (1.25,1.694);
  \node[left,font=\footnotesize]  at (-0.46,1.48) {uniform};
  \node at (0.42,0.12) {(b) Lorentz};
\end{scope}

\end{tikzpicture}
\caption{Where the $4/3$ goes. Both panels show the same accelerating world
line $\gamma$ and the same slab of it, and differ only in what is being called
``at the same moment''. \emph{(a)} Dirac's tube is built from the particle's own
succession of rest frames. Because the velocity changes along $\gamma$,
successive slices are not parallel: they fan open on the side the particle is
accelerating towards, and the slab is wider there. That fanning is the factor
$1+\varepsilon\,a\cdot n$ of Eq.~(\ref{eq:tubeelement}). \emph{(b)} Collecting
field momentum at constant time in one chosen frame gives parallel slices and a
slab of uniform thickness, with no such factor. The wedge by which the two
shaded regions differ is worth exactly one third of the electromagnetic mass,
which is the whole of the discrepancy between Eqs.~(\ref{eq:massren}) and
(\ref{eq:divterm}).
\label{fig:simultaneity}}
\end{figure}
